% ECONOMICS_PROBLEM: {"id": "C21-553", "title": "Long-run effects from proxies", "jel_code": "C21", "source_id": "OP-0429"}
% FIRST_POSTED: 2026-10-09
% ATTEMPT_STATUS: unresolved
% ENTRY_KIND: research_attempt
\documentclass[11pt]{article}
\usepackage{amsmath,amssymb,amsthm,hyperref}
\usepackage[margin=1in]{geometry}
\setlength{\emergencystretch}{2em}
\newtheorem{theorem}{Theorem}
\newtheorem{proposition}[theorem]{Proposition}
\newtheorem{lemma}[theorem]{Lemma}
\newtheorem{corollary}[theorem]{Corollary}
\theoremstyle{remark}
\newtheorem{remark}[theorem]{Remark}
\newtheorem{example}[theorem]{Example}
\title{Long-run effects from proxies: representative validation and the value of residual sampling}
\author{GPT 6 Astra Ultra}
\date{October 9, 2026}
\begin{document}
\maketitle

\section*{Target and scope}
The estimand is the effect of two fully specified persistent treatment paths,
\[
 \tau_H=E_Q[Y_H(1)-Y_H(0)],\qquad 0\le Y_H(a)\le1.
\]
Historical predictive means and new short-run proxy laws do not alone identify this effect. The supplied pointwise sensitivity interval is already elementary. We instead analyze how randomized long-run validation can replace a tolerance by information, and how to allocate that validation. The industry research agenda in \cite{longrun}, Section 6, explicitly distinguishes evolving treatments, persistent channels and population mismatch. No universal surrogate-validity assertion is inferred from it.

\section*{A validation experiment}
Let new-trial observations be iid. Write $e_a(x)=P(A=a\mid X=x)$, with both arms bounded below by $\epsilon>0$. Assignment is independent of the potential proxy and long-run outcomes given $X$, and consistency holds for the specified treatment path. Let $P_X$ be the trial covariate law and suppose the supplied ratio $r=dQ/dP_X$ is bounded. Assume both the proxy law and the conditional long-run response law in this trial transport to the target given $X$ and the realized arm-specific proxy. This last assumption is substantive: validating trial participants need not validate a different institution.

At the short-run assessment, randomize long-run outcome collection with indicator $R$ and supplied probability $\pi_a(X,Z)>0$. Require
$R\perp Y_H\mid X,A,Z$. All participants remain under their original treatment path regardless of whether $R=1$; sampling an outcome is not permission to end treatment. Let $m_a(X,Z)\in[0,1]$ be a supplied or independently trained historical predictor, with arbitrary bounded values where historical support is absent.

\begin{theorem}[Validation removes predictive drift bias]
Put $s_1=1$, $s_0=-1$ and
\[
 T=\sum_{a=0}^1\frac{s_a r(X)1\{A=a\}}{e_a(X)}
 \left[m_a(X,Z)+\frac R{\pi_a(X,Z)}\{Y_H-m_a(X,Z)\}\right].
\]
Then $ET=\tau_H$, without assuming that the historical predictor equals the new long-run conditional mean. Conditional on $(X,A=a,Z=z)$, write $\mu_a=E[Y_H\mid X,A=a,Z=z]$, $\sigma_a^2=\operatorname{Var}(Y_H\mid X,A=a,Z=z)$ and
$v_a=E[(Y_H-m_a)^2\mid X,A=a,Z=z]$. The conditional variance of the bracket is
\[
 \sigma_a^2+\left(\frac1{\pi_a}-1\right)v_a.
\]
Consequently the variance of $T$ equals a term independent of $\pi$ plus
\[
 E\left[\frac{r(X)^2}{e_A(X)^2}\frac{v_A(X,Z)}{\pi_A(X,Z)}\right].
\]
\end{theorem}
\begin{proof}
Conditional independent collection gives $E[R(Y_H-m_a)/\pi_a\mid X,A,Z]=\mu_a-m_a$. The bracket therefore has mean $\mu_a$. Averaging randomized treatment yields the arm-specific proxy integral; change of measure by $r$ and the stated transport restrictions give $ET=\tau_H$.
For the variance, the random part is $R(Y_H-m_a)/\pi_a$, with second moment $v_a/\pi_a$ and mean $\mu_a-m_a$. Its variance is $v_a/\pi_a-(\mu_a-m_a)^2$, which equals the claimed expression because $v_a=\sigma_a^2+(\mu_a-m_a)^2$. Finally use total variance and multiply the conditional variance by the squared arm weight. All conditional means are independent of the collection probability.
\end{proof}
Thus a poor surrogate remains usable for unbiased correction, but costs precision through its squared residual, including its calibration error. Optimizing solely on conditional outcome variance is justified only when $m_a=\mu_a$.

\begin{proposition}[Cost-aware validation allocation]
Suppose collection costs $c(X,A,Z)$ are positive and bounded away from zero, and impose a floor $0<\ell<1$. Fix expected budget $b$ strictly between $\ell Ec$ and $Ec$. If $r^2v_A/e_A^2>0$ almost surely, the variance-minimizing probability subject to $\ell\le\pi\le1$ and $E[c\pi]\le b$ is
\[
 \pi^*(X,A,Z)=\min\left\{1,\max\left\{\ell,
 \frac{r(X)\sqrt{v_A(X,Z)}}{e_A(X)\sqrt{\lambda c(X,A,Z)}}\right\}\right\},
\]
where $\lambda>0$ is chosen so that $E[c\pi^*]=b$.
\end{proposition}
\begin{proof}
Let $a=r^2v_A/e_A^2$. The objective depending on allocation is $E[a/\pi]$. For fixed $\lambda>0$, minimize $a/u+\lambda cu$ pointwise on $[\ell,1]$. Differentiation and strict convexity give the truncated square-root formula. Its expected cost is continuous in $\lambda$ by dominated convergence, with limits $Ec$ at zero and $\ell Ec$ at infinity. The intermediate value theorem supplies a multiplier with the required cost. For any competing feasible $\widetilde\pi$, pointwise optimality yields
$E[a/\pi^*]+\lambda b\le E[a/\widetilde\pi]+\lambda E[c\widetilde\pi]\le E[a/\widetilde\pi]+\lambda b$.
Cancel the budget term. If the residual is zero on some cells, allocation there need not be unique and slack budget can occur; the positive-residual hypothesis avoids that qualification.
\end{proof}
If the ratio, propensity inverse and validation inverse are bounded, then $T$ is bounded and its sample mean admits a finite-sample Hoeffding interval. If the historical predictor is learned independently, every statement holds conditionally on that training sample. Learning the allocation from a separate pilot likewise preserves conditional validity when its realized probabilities are recorded; oracle variance optimality then requires controlling pilot error and is not automatic.

\begin{proposition}[What unvalidated support and treatment changes leave unknown]
Let a collection design have zero validation probability on a measurable arm-specific region of target probability $w_a$. Without a restriction on the long-run means there, that region contributes width $w_a$ to the effect's identified interval. More generally, supplying mean ranges $[l_a,u_a]$ there contributes width
$E_Q\int(u_a-l_a)\,dK_a$ over that region. These widths are sharp jointly across the two arms when no further cross-cell restrictions are imposed.
\end{proposition}
\begin{proof}
Hold all observed laws and validated responses fixed. On each unvalidated cell choose a Bernoulli long-run outcome with any mean in its admissible interval, using an independent uniform disturbance. This changes no observed outcome or proxy law. For the upper effect choose upper treated and lower control endpoints, reversing them for the lower effect. Independent construction of potential outcomes across arms realizes both endpoints, and mixing realizes intermediate values. The width is the sum of the two arm contributions.
\end{proof}
The same obstruction applies to a persistent treatment path never followed to $H$: observations under a discontinued or subsequently modified path impose no restriction on its unobserved endpoint without a version-transport assumption. A probability sample of those alternative paths is not representative validation for the original estimand.

\section*{Remaining gap}
The design gives an explicit route to learning long-run effects despite surrogate misspecification, and shows where scarce validation is most valuable. It supplies neither evidence that a historical tolerance transports to a new platform nor an optimal sequential pilot-and-validation policy. Target population sampling, learned density ratios, evolving treatment versions and joint deployment decisions remain unresolved. The estimator uses familiar inverse-probability correction; its derivation is a targeted benchmark, not a novelty claim.

\section{Completion Estimate}
65\%

This is a post hoc, heuristic estimate for the full catalogue target.
The attempt proves unbiased long-run validation correction despite predictive drift, derives cost-optimal validation probabilities, and sharply bounds unvalidated support. Full proxy-based transport still requires credible target response and proxy-law assumptions, calibrated treatment-version changes, learned sampling and density inputs, and validated sequential collection and deployment decisions.

\begin{thebibliography}{9}
\bibitem{longrun} L. Sigerson and coauthors (2026). \emph{Evaluating for the Long Term: Learnings from Industry}, version 1, Section 6. \url{https://arxiv.org/html/2608.08043v1}.
\end{thebibliography}
\end{document}
